% 01-introducao.tex -- incluído por main-pt.tex
% =====================================================================
\section{Introdução}
\label{sec:intro}
% =====================================================================

A tecnologia blockchain tem recebido grande atenção nos últimos anos, uma vez que pode tratar as questões de confiança, rastreabilidade e dupla contagem nos mercados de carbono florestal sem depender necessariamente de intermediários centralizados \citep{Howson2019ncc,Kotsialou2022}. Em razão disso, diversas contribuições investigaram o uso de registros distribuídos na governança dos mercados de carbono e na gestão dos créditos gerados por projetos de Redução de Emissões por Desmatamento e Degradação Florestal (REDD+) (ver, por exemplo, \citep{Howson2019geoforum,Kotsialou2021,Marke2022,Vilkov2023,He2022} e as referências neles contidas). No Brasil, essa discussão ganhou uma dimensão institucional com a Lei nº 15.042/2024, que instituiu o Sistema Brasileiro de Comércio de Emissões (SBCE), com um registro central administrado por um órgão gestor, obrigações periódicas de relato e de conciliação para os operadores regulados e mecanismos explícitos contra a dupla contagem de créditos jurisdicionais \citep{Lei15042}. Torna-se, assim, concreta a pergunta sobre se um registro em escala nacional poderia ser sustentado por consenso bizantino.

As avaliações quantitativas disponíveis não respondem a essa pergunta. Os trabalhos de \citet{Caburlao2025reic} e \citet{Caburlao2026clei} compararam, por meio do simulador JABS \citep{Yajam2023}, o tempo e o tráfego de finalização de bloco de três projetos REDD+ em operação, Ambify \citep{Ambipar2021}, MCO2 \citep{Moss2020} e TreeCycle \citep{GlobalTree2019}, e mostraram que esses indicadores são determinados, em grande medida, pelo mecanismo de consenso subjacente. Nesses estudos, o número de nós foi dimensionado pela área florestal monitorada, à razão de um nó por $1.000$ hectares, o que é adequado para comparar projetos individuais. Quando extrapolada para a escala nacional, contudo, essa parametrização conduz a conjuntos de validadores com cerca de $10^{5}$ nós, para os quais nenhum protocolo bizantino conhecido é viável. Além disso, na versão do simulador empregada no desenvolvimento da arquitetura Arq-REDD+ \chk{citar a referência da Arq-REDD+ quando publicada}, o tráfego de consenso do pBFT é obtido por uma estimativa em forma fechada proporcional ao número de nós, e não por medição \chk{confirmar se resultados divulgados anteriormente utilizaram essa estimativa}. Uma estimativa linear não pode revelar se a sobrecarga de comunicação se torna a restrição vinculante, uma vez que, no caso normal do pBFT, o número de mensagens por bloco cresce com o quadrado do número de nós \citep{Castro1999}. Portanto, a viabilidade em escala nacional permanece uma questão em aberto, que não pode ser respondida por uma comparação adicional entre algoritmos.

O presente trabalho parte de uma observação simples, mas com consequências diretas para o projeto do registro: em um registro regulado, o número de validadores não é uma grandeza territorial, e sim uma decisão de governança, limitada pelo número de instituições responsabilizáveis (órgão gestor, entes federativos, verificadores credenciados); a demanda de transações, por sua vez, é uma grandeza de mercado, determinada pelo número de operadores e pela estrutura temporal das obrigações de conciliação. As duas escalas afetam o consenso de maneiras distintas. Enquanto o custo de comunicação de cada rodada de consenso cresce com o número de validadores, o volume de transações é agregado em blocos e afeta sobretudo o tamanho dos blocos e a ocupação dos enlaces. Até onde é do conhecimento dos autores, a viabilidade de registros de carbono baseados em blockchain não foi analisada separando essas duas escalas, nem com base em uma demanda derivada da estrutura do mercado regulado.

Motivado pelo exposto, este trabalho investiga em que condições cada família de protocolos de consenso permanece viável em um registro nacional e qual restrição se esgota primeiro. A hipótese de trabalho, a ser confirmada ou refutada pelos experimentos, é que a viabilidade é mais sensível à escala de governança do que à escala de mercado. As principais contribuições são as seguintes:

\begin{itemize}
\item Formula-se o dimensionamento do registro em duas escalas independentes, a de governança ($n$ validadores) e a de mercado (demanda $\lambda$), e quantifica-se, com tráfego medido em nível de pacote, o efeito de acoplar $n$ à área florestal.
\item Deriva-se um modelo de demanda a partir da estrutura de conciliação periódica do SBCE, com forma fechada para os parâmetros de um processo de Poisson modulado por Markov, em vez de assumir chegadas de Poisson estacionárias.
\item Selecionam-se os protocolos avaliados a partir de requisitos do domínio (finalidade determinística, validadores responsabilizáveis, tolerância bizantina e baixo custo energético), explicitando o papel de cada protocolo retido.
\item Estimam-se elasticidades do tráfego, da latência e da ocupação dos enlaces em relação às duas escalas, com validadores bizantinos, por metamodelos validados fora da faixa de ajuste, e derivam-se fronteiras de viabilidade $\nmax(\lambda)$ que identificam a restrição vinculante em cada regime.
\end{itemize}

O restante do artigo está organizado da seguinte forma: a Seção~\ref{sec:background} apresenta a fundamentação teórica e a Seção~\ref{sec:related} discute os trabalhos relacionados. A Seção~\ref{sec:model} define os requisitos do registro, as duas escalas, o modelo de demanda e os critérios de viabilidade. A Seção~\ref{sec:method} detalha o projeto experimental e o procedimento de análise. Resultados numéricos e discussões são apresentados na Seção~\ref{sec:results}. Por fim, as considerações finais são apresentadas na Seção~\ref{sec:conclusion}.
